% GENERATED FILE. Edit canonical lessons, then run npm run generate:books.
% canonical-source-hash: fnv1a64:c56a970efd5965ed
% canonical-lessons: JA-C387-demukaeru, JA-C387-mikakeru, JA-C387-toujou, JA-C387-koukuuken, JA-C387-shutchou

\chapter{Meeting the Plane}
\label{ch:ja-meeting-the-plane}
\hlchaptermodality{387}

\begin{chapteropening}
\textbf{By the end of this chapter:} \emph{I can meet someone arriving at the airport.}

Complete the last lesson of chapter 387: I can meet someone arriving at the airport.
\end{chapteropening}

\section[\texorpdfstring{\ja{でむかえる} (demukaeru)}{でむかえる (demukaeru)}]{\ja{でむかえる} (demukaeru) — to go out to meet (someone arriving)}

\label{lesson:JA-C387-demukaeru}



\begin{quote}
Before the new one: say the Japanese for departure from a country, then the Japanese for to take something out.
\end{quote}

\subsection*{You'll want to know: \ja{でむかえる}}
\textbf{\ja{でむかえる}} — \emph{demukaeru} — \textquotedblleft{}to go out to meet (someone arriving)\textquotedblright{}.

\begin{quote}
\textbf{\ja{くうこうで} \ja{そふを} \ja{でむかえました。}} — \emph{kūkō de sofu o demukaemashita} — \textquotedblleft{}We met my grandfather at the airport.\textquotedblright{}
\end{quote}

Polite \textbf{\ja{でむかえます}}, \emph{demukaemasu}; an ichidan verb. Past (polite): \textbf{\ja{ははが} \ja{くうこうで} \ja{そふを} \ja{でむかえました。}}, \emph{haha ga kūkō de sofu o demukaemashita}.

\begin{cousinweb}
\textbf{\ja{でむかえる}} (\emph{demukaeru}) is a native Japanese word (wago), inherited rather than borrowed from Chinese, built from \emph{de}, \textquotedblleft{}going out\textquotedblright{} (from \emph{deru}) and \emph{mukaeru}, \textquotedblleft{}to welcome\textquotedblright{}; its kanji were attached later as kun'yomi (JMdict).
\end{cousinweb}

\begin{grammarlens}[title={What you've built}]
One more word for getting around.
\end{grammarlens}

\subsection*{Your turn}
\begin{itemize}
\raggedright
  \item \emph{Say it:} \emph{demukaeru}
  \item \emph{Say it:} \emph{demukaeru}, once more
  \item \emph{Say it:} \emph{kūkō de sofu o demukaemashita} — We met my grandfather at the airport.
  \item \emph{Recall:} say the Japanese for departure from a country, then the Japanese for to take something out, then say \emph{demukaeru} again
\end{itemize}

\begin{tcolorbox}[breakable,colback=teal!4,colframe=teal!35!black,title={Before you move on}]
What does \ja{でむかえる} mean? (\textquotedblleft{}To go out to meet (someone arriving)\textquotedblright{}.) Say it once more.
\end{tcolorbox}
\section[\texorpdfstring{\ja{みかける} (mikakeru)}{みかける (mikakeru)}]{\ja{みかける} (mikakeru) — to happen to see; to catch sight of}

\label{lesson:JA-C387-mikakeru}



\begin{quote}
Before the new one: say the Japanese for to take something out, then the Japanese for to go out to meet.
\end{quote}

\subsection*{You'll want to know: \ja{みかける}}
\textbf{\ja{みかける}} — \emph{mikakeru} — \textquotedblleft{}to happen to see; to catch sight of\textquotedblright{}.

\begin{quote}
\textbf{\ja{くうこうで} \ja{せんせいを} \ja{みかけました。}} — \emph{kūkō de sensei o mikakemashita} — \textquotedblleft{}I happened to see my teacher at the airport.\textquotedblright{}
\end{quote}

Polite \textbf{\ja{みかけます}}, \emph{mikakemasu}; an ichidan verb. Past (polite): \textbf{\ja{あねが} \ja{えきで} \ja{せんせいを} \ja{みかけました。}}, \emph{ane ga eki de sensei o mikakemashita}.

\begin{cousinweb}
\textbf{\ja{みかける}} (\emph{mikakeru}) is a native Japanese word (wago), inherited rather than borrowed from Chinese; its kanji was attached later as a kun'yomi (JMdict). Its first part is \emph{mi}, from \emph{miru}, \textquotedblleft{}to see\textquotedblright{} (JMdict).
\end{cousinweb}

\begin{grammarlens}[title={What you've built}]
One more word for getting around.
\end{grammarlens}

\subsection*{Your turn}
\begin{itemize}
\raggedright
  \item \emph{Say it:} \emph{mikakeru}
  \item \emph{Say it:} \emph{mikakeru}, once more
  \item \emph{Say it:} \emph{kūkō de sensei o mikakemashita} — I happened to see my teacher at the airport.
  \item \emph{Recall:} say the Japanese for to take something out, then the Japanese for to go out to meet, then say \emph{mikakeru} again
\end{itemize}

\begin{tcolorbox}[breakable,colback=teal!4,colframe=teal!35!black,title={Before you move on}]
What does \ja{みかける} mean? (\textquotedblleft{}To happen to see; to catch sight of\textquotedblright{}.) Say it once more.
\end{tcolorbox}
\section[\texorpdfstring{\ja{とうじょう} (tōjō)}{とうじょう (tōjō)}]{\ja{とうじょう} (tōjō) — boarding (a plane or ship)}

\label{lesson:JA-C387-toujou}



\begin{quote}
Before the new one: say the Japanese for to go out to meet, then the Japanese for to happen to see.
\end{quote}

\subsection*{You'll want to know: \ja{とうじょう}}
\textbf{\ja{とうじょう}} — \emph{tōjō} — \textquotedblleft{}boarding (a plane or ship)\textquotedblright{}.

\begin{quote}
\textbf{\ja{とうじょうの} \ja{じかんは} \ja{くじです。}} — \emph{tōjō no jikan wa kuji desu} — \textquotedblleft{}Boarding time is nine o'clock.\textquotedblright{}
\end{quote}

\begin{cousinweb}
\textbf{\ja{とうじょう}} (\emph{tōjō}) is Sino-Japanese (kango): two Chinese characters read with their on'yomi, readings borrowed from Chinese. They are \emph{tō}, \textquotedblleft{}board\textquotedblright{} (Middle Chinese \emph{thap}) and \emph{jō}, \textquotedblleft{}ride\textquotedblright{} (\emph{zying}) (KANJIDIC2; Middle Chinese in Baxter's transcription).
\end{cousinweb}

\begin{grammarlens}[title={What you've built}]
One more word for getting around.
\end{grammarlens}

\subsection*{Your turn}
\begin{itemize}
\raggedright
  \item \emph{Say it:} \emph{tōjō}
  \item \emph{Say it:} \emph{tōjō}, once more
  \item \emph{Say it:} \emph{tōjō no jikan wa kuji desu} — Boarding time is nine o'clock.
  \item \emph{Recall:} say the Japanese for to go out to meet, then the Japanese for to happen to see, then say \emph{tōjō} again
\end{itemize}

\begin{tcolorbox}[breakable,colback=teal!4,colframe=teal!35!black,title={Before you move on}]
What does \ja{とうじょう} mean? (\textquotedblleft{}Boarding (a plane or ship)\textquotedblright{}.) Say it once more.
\end{tcolorbox}
\section[\texorpdfstring{\ja{こうくうけん} (kōkūken)}{こうくうけん (kōkūken)}]{\ja{こうくうけん} (kōkūken) — a plane ticket}

\label{lesson:JA-C387-koukuuken}



\begin{quote}
Before the new one: say the Japanese for to happen to see, then the Japanese for boarding (a plane or ship).
\end{quote}

\subsection*{You'll want to know: \ja{こうくうけん}}
\textbf{\ja{こうくうけん}} — \emph{kōkūken} — \textquotedblleft{}a plane ticket\textquotedblright{}.

\begin{quote}
\textbf{\ja{なつやすみの} \ja{こうくうけんは} \ja{たかかったです。}} — \emph{natsuyasumi no kōkūken wa takakatta desu} — \textquotedblleft{}The plane tickets for the summer holidays were expensive.\textquotedblright{}
\end{quote}

\begin{cousinweb}
\textbf{\ja{こうくうけん}} (\emph{kōkūken}) is Sino-Japanese (kango): three Chinese characters read with their on'yomi, readings borrowed from Chinese. They are \emph{kō}, \textquotedblleft{}navigate\textquotedblright{} (Middle Chinese \emph{hang}), \emph{kū}, \textquotedblleft{}sky\textquotedblright{} (\emph{khuwng}) and \emph{ken}, \textquotedblleft{}ticket\textquotedblright{} (\emph{khjwonH}) (KANJIDIC2; Middle Chinese in Baxter's transcription).
\end{cousinweb}

\begin{grammarlens}[title={What you've built}]
One more word for getting around.
\end{grammarlens}

\subsection*{Your turn}
\begin{itemize}
\raggedright
  \item \emph{Say it:} \emph{kōkūken}
  \item \emph{Say it:} \emph{kōkūken}, once more
  \item \emph{Say it:} \emph{natsuyasumi no kōkūken wa takakatta desu} — The plane tickets for the summer holidays were expensive.
  \item \emph{Recall:} say the Japanese for to happen to see, then the Japanese for boarding (a plane or ship), then say \emph{kōkūken} again
\end{itemize}

\begin{tcolorbox}[breakable,colback=teal!4,colframe=teal!35!black,title={Before you move on}]
What does \ja{こうくうけん} mean? (\textquotedblleft{}A plane ticket\textquotedblright{}.) Say it once more.
\end{tcolorbox}
\section[\texorpdfstring{\ja{しゅっちょう} (shutchō)}{しゅっちょう (shutchō)}]{\ja{しゅっちょう} (shutchō) — a business trip}

\label{lesson:JA-C387-shutchou}



\begin{quote}
Before the new one: say the Japanese for boarding (a plane or ship), then the Japanese for a plane ticket.
\end{quote}

\subsection*{You'll want to know: \ja{しゅっちょう}}
\textbf{\ja{しゅっちょう}} — \emph{shutchō} — \textquotedblleft{}a business trip\textquotedblright{}.

\begin{quote}
\textbf{\ja{ちちは} \ja{しゅっちょうで} \ja{がいこくに} \ja{いきました。}} — \emph{chichi wa shutchō de gaikoku ni ikimashita} — \textquotedblleft{}My father went abroad on a business trip.\textquotedblright{}
\end{quote}

\begin{cousinweb}
\textbf{\ja{しゅっちょう}} (\emph{shutchō}) is Sino-Japanese (kango): two Chinese characters read with their on'yomi, readings borrowed from Chinese. They are \emph{shutsu}, \textquotedblleft{}exit\textquotedblright{} (Middle Chinese \emph{tsyhwit}) and \emph{chō}, \textquotedblleft{}stretch\textquotedblright{} (\emph{trjang}), joined as \emph{shutchō} (KANJIDIC2; Middle Chinese in Baxter's transcription).
\end{cousinweb}

\begin{grammarlens}[title={What you've built}]
One more word for getting around.
\end{grammarlens}

\subsection*{Your turn}
\begin{itemize}
\raggedright
  \item \emph{Say it:} \emph{shutchō}
  \item \emph{Say it:} \emph{shutchō}, once more
  \item \emph{Say it:} \emph{chichi wa shutchō de gaikoku ni ikimashita} — My father went abroad on a business trip.
  \item \emph{Recall:} say the Japanese for boarding (a plane or ship), then the Japanese for a plane ticket, then say \emph{shutchō} again
\end{itemize}

\begin{tcolorbox}[breakable,colback=teal!4,colframe=teal!35!black,title={Before you move on}]
What does \ja{しゅっちょう} mean? (\textquotedblleft{}A business trip\textquotedblright{}.) Say it once more.
\end{tcolorbox}
